% Generated draft: requires mathematical and source review.
\documentclass[11pt]{article}
\usepackage[margin=25mm]{geometry}
\usepackage{fontspec}
\setmainfont{Latin Modern Roman}
\usepackage{amsmath,amssymb,mathtools,unicode-math}
\setmathfont{Latin Modern Math}
\usepackage{enumitem,needspace,xurl,hyperref}
\setlength{\parindent}{0pt}
\setlength{\parskip}{5pt}
\setlength{\emergencystretch}{3em}
\newcommand{\N}{\mathbb{N}}
\newcommand{\Z}{\mathbb{Z}}
\newcommand{\Q}{\mathbb{Q}}
\newcommand{\R}{\mathbb{R}}
\newcommand{\C}{\mathbb{C}}
\newcommand{\F}{\mathbb{F}}
\newcommand{\A}{\mathbb{A}}
\newcommand{\PP}{\mathbb{P}}
\newcommand{\E}{\mathbb{E}}
\newcommand{\eps}{\varepsilon}
\newcommand{\dd}{\,\mathrm{d}}
\begin{document}
\section*{UnsolvedMath ID: 2302063}
\textbf{Source problem code:} catalogue entry
\subsection{Definitions and mathematical statement}
Let $f$ be a rational map of the Riemann sphere $\widehat\C$, and let $f^n$ denote its $n$-fold iterate, not its pointwise $n$th power. The set $C$ is specified only by a cross-reference to Problem 2.62, absent from this record. A customary intended choice is the Fatou set, where the family of iterates is normal; this identification requires checking the missing source context.

On a connected component $G$ of the intended $C$, a limit function is a subsequential limit
\[
 f^{n_j}|_G\longrightarrow g\quad(n_j\to\infty),
\]
with local uniform convergence in the spherical metric under the usual dynamical convention. A constant limit has value $\alpha\in\widehat\C$. When counting constant limits across different components, the natural value-based set is
\[
 A(f)=\{\alpha:\exists G\subset C\text{ a component and }n_j\to\infty,
                  \ f^{n_j}|_G\to\alpha\}.
\]
The source does not clarify whether it counts values or domain-function pairs.

The source says that in its simplest examples the limit functions form a finite set and are periodic constants: $f^k(\alpha)=\alpha$ for some $k\geq1$. It then defines class $N$ by additionally requiring
\[
 |(f^k)'(\alpha)|<1
\]
for each constant limit. At infinity the multiplier is evaluated in a local coordinate. The wording appears to include the preceding finiteness/constant-limit assumptions; whether they are part of the formal definition of $N$ must be checked against the full source.

\textbf{Questions (I. N. Baker).} (1) Does there exist a rational $f$ with infinitely many constant limit functions, in particular $|A(f)|=\infty$ under the value-counting convention? (2) Is membership in $N$ generic in some sense among rational functions?

The source does not prescribe a degree range, a fixed-degree parameter space, a topology, or whether generic means residual, open dense, full measure, or something else. These are essential choices for the second question; no density or hyperbolicity conjecture is substituted.

The supplied research summary labels the item PARTIAL-PROGRESS. It reports that Eremenko--Lyubich obtained infinitely many constant limits for transcendental entire functions, referring to Problem 2.67, but this is a different function class from the rational maps asked about here. It mentions links with Sullivan/McMullen rigidity, says generic properties are not fully settled, and reports no complete resolution located. These are unverified supplied status claims.
\subsection{Short English statement}
Can a rational dynamical system have infinitely many constant limits of subsequences of its iterates on normality components? Is the class whose limit functions are finitely many attracting periodic constants generic? The referenced set $C$, exact class-$N$ definition, degree range, and meaning of generic are not supplied.
\subsection{Sources}
\begin{itemize}
\item Local UnsolvedMath registry, permanent ID 2302063, source code catalogue entry, ``Research Problems in Function Theory --- Problem 2.63''; attribution: I. N. Baker. Supplied status partially solved, classification PARTIAL-PROGRESS, updated 2026-08-06.
\item Hayman, \emph{Research Problems in Function Theory} (2018), Problem 2.63, with source cross-references to Problems 2.62 and 2.67, \url{https://arxiv.org/abs/1809.07200}. The record reports source-TeX extraction and access on 2026-07-29; the 2018 update was retained for status triage. Current status and rights require review; the record describes publicly available arXiv source TeX.
\item The statement mentions older accounts by Fatou with empty citation slots and a modern sketch by Guckenheimer without a complete reference. The research summary mentions Eremenko--Lyubich and Sullivan/McMullen without particular publications. No missing bibliographic details are invented.
\end{itemize}
\end{document}
